\chapter{Heuristic Transformations and Working Backwards}
\label{sec:p12-heuristic-transformations-and-backward-reasoning}

A useful intermediate problem must be connected to the problem that
was actually posed. The connection may preserve all solutions, enlarge
the candidate set, restrict it, or introduce auxiliary objects from
which a solution can be reconstructed. Recording that connection makes
it possible to explore freely while keeping a rigorous final argument.

P\'olya organizes problem solving around understanding, planning,
execution and reconsideration. His discussions of decomposition,
generalization, specialization and working backwards supply the examples
below \cite[Part~I, secs.~6--14; Part~III, the corresponding dictionary
entries]{Polya2004HowToSolveItSelected}. The formal solution-set and
transition-relation records used here make the logical obligations of
those moves explicit. They are an editorial framework for checking
reasoning, not a claim that P\'olya supplied a complete discovery algorithm.

The mathematical examples use the real numbers and ordinary Euclidean
geometry. For vectors, \(u\cdot v\) is the Euclidean dot product and
\(\|u\|=\sqrt{u\cdot u}\). Angles are interior angles in radians.
Area and volume have their usual Euclidean meanings, with congruence
invariance and computation by parallel cross-sections. The explicit
calculations below identify where these geometric premises enter.

\section{State the problem before changing it}

\begin{hbdefinition}{def:p12-heuristic-problem}{Unknown, data, condition and requested output}
A \emph{problem to find} specifies fixed data \(d\), a domain
\(X_d\) of possible unknowns, and a condition \(C_d(x)\).
Its solution set is
\[
             S_d=\{x\in X_d:C_d(x)\}.
\]
The request must also specify whether to find one solution, all
solutions, a solution optimizing a stated objective, or a construction
using specified operations. An existence claim is
\(\exists x\in X_d\,C_d(x)\); it need not provide a construction
in the stipulated system of operations.

A \emph{problem to prove} specifies assumptions \(\Gamma\) and a
target assertion \(T\), with an intended meaning and domain. Its
deductive task is to establish \(\Gamma\vdash T\) in the stated
proof system. If the request is to decide whether the assertion follows,
a countermodel satisfying \(\Gamma\) and falsifying \(T\) is a
possible conclusive answer. The assumptions may be empty: a theorem
need not be written as a visible if--then sentence.

Record auxiliary assumptions separately from \(\Gamma\), and auxiliary
unknowns separately from the required output. A proposed diagram is a
representation of these data and conditions; its apparent symmetry,
perpendicularity or scale is not an additional premise.
\index{problem solving!problem specification}
\index{unknown!data and conditions}
\end{hbdefinition}

\begin{hbexample}{def:p12-heuristic-types}{The same data can express different problems}
Three positive lengths can prescribe a triangle's sides or its three
altitudes. Both requests seek a triangle and have three numerical data,
but their conditions differ. The lengths \(1,1,3\) cannot be the
sides of a nondegenerate Euclidean triangle, yet they can be its
altitudes; \cref{thm:p12-heuristic-geometric-specifications} supplies
the construction and verification.

A word puzzle seeks a word satisfying both semantic and spelling
conditions. A chess problem seeks a legal move with a specified effect,
not merely a visually attractive displacement. A claim about a witness's
account requires an evidential standard appropriate to that inquiry.
Calling these all problems does not make their output or evidence
criteria interchangeable. P\'olya's distinction between finding an object
and establishing an assertion directs attention to that difference
\cite[Part~III, ``Problems to find, problems to prove'' and
``What is the unknown?'']{Polya2004HowToSolveItSelected}.
\end{hbexample}

\begin{hbmethod}{def:p12-heuristic-four-phases}{Four recurring tasks}
\begin{enumerate}
\item \emph{Understand.} State the requested object or assertion, its
domain, all data and every condition. Identify ambiguities and determine
which observations are assumptions and which are only guesses.
\item \emph{Plan.} Search for a relevant solved problem, theorem,
representation, auxiliary construction or transformed task. Write the
logical connection that would make the proposed plan useful.
\item \emph{Execute.} Perform the chosen calculations or constructions
and justify their side conditions. Keep rejected cases and unresolved
obligations visible until they have been settled.
\item \emph{Reconsider.} Check the answer against the original condition,
examine vulnerable steps, and ask whether another derivation, a special
case, symmetry or a dimensional calculation exposes a mistake or a
further application.
\end{enumerate}
These tasks can recur: a failed calculation can send the solver back
to the interpretation or plan. Their order is a guide to organizing
work, not a fixed number of mental steps or an empirical guarantee of
success \cite[Part~I, secs.~6--14]{Polya2004HowToSolveItSelected}.
\end{hbmethod}

\section{Decompose conditions and preserve the direction of inference}

\begin{hbdefinition}{def:p12-heuristic-transformation}{Three relations between a target and a transformed condition}
Fix the original data and domain, and abbreviate the condition by
\(C(x)\). For any condition \(B\) on that fixed domain, write
\(S_B=\{x\in X_d:B(x)\}\). The following relations are pointwise for every admissible
\(x\):
\[
\begin{array}{c|c|l}
\text{transformed condition}&\text{relation proved}&\text{permitted use}\\ \hline
N(x)&C(x)\Rightarrow N(x)&\text{generate necessary-condition candidates}\\
H(x)&H(x)\Rightarrow C(x)&\text{find a sufficient route to a solution}\\
E(x)&C(x)\Longleftrightarrow E(x)&\text{replace the condition equivalently}
\end{array}
\]
Writing an implication is an obligation to justify it. It is not
evidence that the implication holds. Domain restrictions are included
in the condition, so a transformation that is valid only for
\(x\ne0\) must retain or separately handle \(x=0\).
\end{hbdefinition}

\begin{hbtheorem}{thm:p12-heuristic-solution-relations}{proved=annex:p12.heuristic.solution-relations}{Candidate sets, decomposition and reconstruction}
The three relations above give, respectively,
\(S_C\subseteq S_N\), \(S_H\subseteq S_C\), and
\(S_E=S_C\). If \(C=\bigwedge_{i=1}^k C_i\), for an integer
\(k\ge0\), then
\[
                     S_C=\bigcap_{i=1}^k S_{C_i}.
\]
Nonemptiness of every individual \(S_{C_i}\) does not imply
nonemptiness of their intersection. For subproblems with shared
variables, their local answers must agree on the shared values and
satisfy any conditions linking the subproblems.

More generally, let an auxiliary problem have solution set \(A\)
and let \(F:A\to X_d\) be a reconstruction map. Proving
\(C_d(F(y))\) for all \(y\in A\) makes every auxiliary solution
a source of an original solution. To obtain \emph{all} original
solutions by this method, one must additionally establish
\(F(A)=S_d\). A map \(G:S_d\to A\) with \(F(G(x))=x\)
is sufficient for that additional condition.
\end{hbtheorem}

P\'olya's recombinations include keeping the unknown while modifying
data or conditions, keeping the data while introducing another unknown,
and changing both. For a proof problem, the corresponding changes concern
hypothesis and conclusion. These choices describe how to generate a
new task; \cref{thm:p12-heuristic-solution-relations} states what must
be checked before its answer can be used
\cite[Part~III, ``Decomposing and recombining'']{Polya2004HowToSolveItSelected}.

For example, specifying a circle separates its center from its radius.
Finding either can be useful without completing the circle problem:
the remaining conditions still have to determine or constrain the other.
Another recombination keeps both the unknown and the data but restates
the condition. Constructing an equilateral triangle on a given side
and constructing an equiangular triangle on that side are equivalent
Euclidean tasks; the angle--side equivalence is included in
\cref{thm:p12-heuristic-geometric-specifications}.

\begin{hbexample}{def:p12-heuristic-word-intersection}{A semantic and spelling intersection}
The source's machine-part puzzle combines a five-letter restriction,
a machine-part meaning, and an interpretation of the clue suggesting
reading in both directions. Candidate words such as \emph{lever},
\emph{screw}, \emph{wheel}, \emph{shaft}, \emph{hinge} and
\emph{motor} satisfy parts of the intended condition. The word
\emph{rotor} also satisfies the reversal requirement: reversing its
five letters leaves the same word. This verifies the proposed answer
under that interpretation of the clue. It does not prove uniqueness
over every possible dictionary or every possible interpretation of a pun.
The proposed interpretation is part of the work that must be checked,
not a spelling identity derived from the original wording alone.
\end{hbexample}

\begin{hbtheorem}{thm:p12-heuristic-triangle-loci}{proved=annex:p12.heuristic.triangle-loci}{A triangle from the intersection of two loci}
Let \(a,h>0\) and \(0<\alpha<\pi\). Seek a Euclidean triangle
with base length \(a\), altitude to that base \(h\), and opposite
angle \(\alpha\). Fix \(B=(0,0)\), \(C=(a,0)\), and choose
the upper half-plane for the third vertex. Then all answers have
\(A=(t,h)\), where
\[
 t=\frac a2\mathbin{\pm}
       \sqrt{\frac{a^2}{4}-h^2+ah\cot\alpha}.
\]
A negative radicand means there is no answer in that half-plane;
zero gives one position, and a positive radicand two positions.
The opposite half-plane gives reflected answers if it is allowed.
The foot of the altitude is not required to lie on the base segment.

Geometrically, one condition gives the line at distance \(h\)
from \(BC\). The angle condition gives the appropriate circular
arc through \(B,C\). Their intersection enforces both conditions
on the \emph{same} vertex. This is the source's two-locus construction
with feasibility and half-plane choices made explicit.
\end{hbtheorem}

\section{Introduce an auxiliary unknown}

\begin{hbtheorem}{thm:p12-heuristic-box}{proved=annex:p12.heuristic.box}{The box diagonal, checks and a cable application}
For a rectangular box with edge lengths \(a,b,c>0\), the length
\(x\) of a space diagonal is
\[
                  x=\sqrt{a^2+b^2+c^2}.
\]
Introducing the face diagonal \(y\) gives the two right-triangle
relations \(y^2=a^2+b^2\), \(x^2=y^2+c^2\). The box center
is the midpoint of each of its four space diagonals, so its distance
to any vertex is \(x/2\).

The expression is symmetric in the three lengths, increases when
one positive length increases, scales by \(\lambda\) when all
lengths scale by \(\lambda>0\), and approaches
\(\sqrt{a^2+b^2}\) as \(c\) tends to zero. These are checks
of the result and its interpretation; they supplement the derivation.

For a \(21\)-by-\(16\)-yard rectangular roof with an \(8\)-yard
pole at its center, a straight cable from a roof corner to a point
\(2\) yards below the pole's top has length
\[
                  \sqrt{10.5^2+8^2+6^2}=14.5\ \text{yards}.
\]
The four corners give equal lengths. The idealized geometry specifies
straight, taut segments; it does not include sag or installation allowance.
\end{hbtheorem}

These are P\'olya's examples in Part~I, Sections~8, 10, 12 and 14.
The face diagonal changes the immediate unknown while retaining the
original measurements. Its value is useful because the second
right triangle reconstructs the requested space diagonal. The center
and roof applications reuse either that result or the two-triangle
method. The four diagonals and the six face-based pyramids meeting
at the center offer further ways to examine the same configuration.

\begin{hbtheorem}{thm:p12-heuristic-geometric-specifications}{proved=annex:p12.heuristic.geometric-specifications}{Checking the exact geometric condition}
There is no nondegenerate Euclidean triangle with side lengths
\(1,1,3\). There is a triangle with altitude lengths \(1,1,3\):
one example has vertices
\[
 (0,3),\qquad(-3/\sqrt{35},0),\qquad(3/\sqrt{35},0).
\]
Also, if a Euclidean quadrilateral has four equal side lengths and
both diagonals are nonzero, its diagonals are perpendicular.
This assertion requires the side-equality hypothesis; it is a
different task from constructing a triangle from given lengths.
For a nondegenerate Euclidean triangle, equality of all three sides
is equivalent to equality of all three interior angles.
\end{hbtheorem}

\section{Work backwards, then justify forwards}

\begin{hbtheorem}{thm:p12-heuristic-backward-proof}{proved=annex:p12.heuristic.backward-proof}{The implication directions needed by backward proof search}
Use a fixed assumption context and a deductive system admitting
modus ponens in that context. Suppose \(T_0\) is the desired assertion.
If each candidate predecessor
is proved sufficient, so that
\[
 \Gamma\vdash T_i\to T_{i-1}\quad(1\le i\le n),
 \qquad\Gamma\vdash T_n,
\]
then \(\Gamma\vdash T_0\). A chain only in the other direction,
\(\Gamma\vdash T_{i-1}\to T_i\), does not license this conclusion.
Such a chain supplies necessary conditions that may be useful for
candidate generation, but the candidate still requires a forward
verification. If every step is an equivalence under its recorded
side conditions, it may be traversed in either direction while
retaining those conditions.
\end{hbtheorem}

\begin{hbtheorem}{thm:p12-heuristic-square-root}{proved=annex:p12.heuristic.square-root}{A necessary algebraic condition and its rejected root}
Over the real numbers, the equation
\[
                         \sqrt{x+1}=x-1
\]
has exactly the solution \(x=3\). The original equation implies
\(x\ge1\). Squaring and rearranging gives the necessary condition
\(x(x-3)=0\), with candidates \(0,3\); the first is rejected.
Equivalently, retaining \(x\ge1\) throughout makes the squaring
step reversible and produces the same unique solution.
This example is an editorial application of the transformation record.
\end{hbtheorem}

\begin{hbdefinition}{def:p12-heuristic-predecessor}{Legal states and candidate predecessors}
For a state space \(V\), a legal one-step relation
\(R\subseteq V\times V\), and a goal set \(G\subseteq V\), put
\[
 \operatorname{Pre}_R(G)=
     \{s\in V:\exists t\in G\ (s,t)\in R\}.
\]
A backward step selects a state and a \emph{legal forward transition}
from it to the current target. It need not be an inverse operation:
emptying a container, for example, has many possible predecessor
states, and its reverse is not permission to fill an arbitrary amount.
\end{hbdefinition}

\begin{hbtheorem}{thm:p12-heuristic-backward-path}{proved=annex:p12.heuristic.backward-path}{Predecessor witnesses reconstruct a forward path}
Define \(G_0=G\) and \(G_{k+1}=\operatorname{Pre}_R(G_k)\).
For every integer \(k\ge0\), \(s\in G_k\) exactly when there is a legal path of length
\(k\) from \(s\) to a state in \(G\). Keeping the witnesses
of the backward steps reconstructs this path in forward order.
The statement concerns path validity and existence, not a guarantee
that a particular strategy will find the useful predecessors.
\end{hbtheorem}

\begin{hbtheorem}{thm:p12-heuristic-pails}{proved=annex:p12.heuristic.pails}{Six quarts using nine- and four-quart containers}
Let a state \((L,S)\) record the amounts in containers of capacities
\(9\) and \(4\). Initially both are empty. Allowed actions are
filling a container from the river, emptying it into the river, or
pouring between containers until the source is empty or the destination
is full. The following legal sequence ends with six quarts in the
larger container and an empty smaller one:
\[
\begin{aligned}
(0,0)&\to(9,0)\to(5,4)\to(5,0)\to(1,4)\\
     &\to(1,0)\to(0,1)\to(9,1)\to(6,4)\to(6,0).
\end{aligned}
\]
Backward discovery identifies \((9,1)\) as useful: the smaller
container has three quarts of spare capacity, so a full larger
container can be reduced to six. The earlier steps obtain the
one-quart amount required by this plan.
\end{hbtheorem}

The pail problem and the discussion of analysis in P\'olya's
``Working backwards'' entry illustrate discovering a sequence in
one order and performing it in the other. The entry also describes
a detour around a three-sided fence. A feasible route can begin
by increasing direct distance to the destination. Its brief comparison
of animal behavior is presented there as an illustration with explicit
reservations; it supplies no controlled species comparison, effect
size or general law of human discovery. Empirical claims about
noticing and transferring analogies have their separate treatment
in \cref{sec:p19-analogy-and-problem-transfer}.

\section{Generalization, specialization and justified transfer}

\begin{hbtheorem}{thm:p12-heuristic-generalization}{proved=annex:p12.heuristic.generalization}{What special cases establish}
Let \(E\subseteq D\). A statement \(\forall x\in D\,P(x)\)
implies \(\forall x\in E\,P(x)\). A point \(e\in E\)
with \(\neg P(e)\) refutes the statement on \(D\).
Truth on \(E\) alone does not imply truth on \(D\).

In particular, for every \(N\ge1\), the polynomial
\(p_N(x)=\prod_{j=1}^{N}(x-j)\) vanishes at every checked
integer \(1,\ldots,N\) but not at \(N+1\).
Finitely many matching examples therefore do not, by themselves,
prove an unrestricted algebraic identity.
\end{hbtheorem}

Generalizing a problem can expose a simpler invariant; specializing it
can reveal a tractable model or a counterexample. An analogy can suggest
either move. Deductive transfer then requires an additional argument:
restriction from a proved general result, a reconstruction map, or an
invariance that takes the original case to the solved special case.
Similarity alone supplies none of these implications. This preserves
the distinction made in P\'olya's ``Heuristic reasoning'' entry between
a provisional suggestion and a completed proof.

\begin{hbtheorem}{thm:p12-heuristic-cubes}{proved=annex:p12.heuristic.cubes}{From a cube-sum pattern to an induction proof}
For every integer \(n\ge0\),
\[
                  \sum_{k=1}^{n}k^3=
                        \left(\frac{n(n+1)}2\right)^2.
\]
The instance \(1+8+27+64=100\), used in the source's
``Generalization'' entry, suggests investigating the larger family.
The induction proof establishes every instance; the initial numerical
pattern does not perform that proof.
\end{hbtheorem}

\begin{hbtheorem}{thm:p12-heuristic-central-bisection}{proved=annex:p12.heuristic.central-bisection}{Generalizing the octahedron problem}
Let a bounded measurable solid in Euclidean three-space have a center
of symmetry \(o\), and let \(\ell\) be a line. Every plane
containing both \(o\) and \(\ell\) bisects the solid's volume.
Such a plane is unique when \(o\notin\ell\); when
\(o\in\ell\), any plane through \(\ell\) suffices.
In particular this solves the source's problem for a regular octahedron.
\end{hbtheorem}

\begin{hbtheorem}{thm:p12-heuristic-frustum}{proved=annex:p12.heuristic.frustum}{Replacing numerical data by parameters}
A square-pyramid frustum with parallel square bases of side lengths
\(a>b\ge0\) and perpendicular height \(h>0\) has volume
\[
                    V=\frac h3(a^2+ab+b^2).
\]
The equal-side case \(a=b>0\) is a prism and has volume \(ha^2\),
the same formula. The source's numerical data \(a=10\), \(b=5\),
\(h=6\) give \(V=350\) cubic inches. The parameterized expression
also makes scaling, the pyramid case \(b=0\), and the prism case
available as checks.
\end{hbtheorem}

\begin{hbtheorem}{thm:p12-heuristic-triangle-counterexample}{proved=annex:p12.heuristic.triangle-counterexample}{A genuine triangle refuting a proposed inequality}
Let \(r\) be a triangle's inradius, \(R\) its circumradius,
and \(H\) its longest altitude. The proposed universal inequality
\(r+R\le H\) holds with equality for equilateral triangles,
but is false in general. For the nondegenerate triangle
\[
                   (-1,0),\quad(1,0),\quad(0,1/4)
\]
one has
\[
 r=\frac1{4+\sqrt{17}},\qquad R=\frac{17}{8},\qquad
 H=\frac2{\sqrt{17}},
\]
and \(r+R>H\). Thus no appeal to a degenerate limiting triangle
is necessary. This is a finite realization of the flattening
specialization suggested by the source.
For an isosceles triangle with its two equal sides fixed at one and
vertex angle \(2\theta\), \(0<\theta<\pi/2\), the formulas are
\[
 r=\frac{\sin\theta\cos\theta}{1+\sin\theta},\qquad
 R=\frac1{2\cos\theta},\qquad
 H=\max\{\cos\theta,2\sin\theta\cos\theta\}.
\]
They justify the source's exploratory limiting cases: as \(\theta\)
tends to zero, \((r,R,H)\) tends to \((0,1/2,1)\); as
\(\theta\) tends to \(\pi/2\), \(r,H\) tend to zero and
\(R\) diverges. The explicit positive-area counterexample supplies
the actual refutation within the problem's domain.
\end{hbtheorem}

\begin{hbtheorem}{thm:p12-heuristic-ships}{proved=annex:p12.heuristic.ships}{Turning one moving ship into a stationary point}
In the idealized point model, let the two positions at time \(t\)
be \(a+tu\) and \(b+tv\) in Euclidean space. Set
\(d=a-b\), \(w=u-v\). Their distance is \(\|d+tw\|\).
If \(w=0\), it is constant. If \(w\ne0\), its minimum over
all real times occurs at
\[
                   t_*=-\frac{d\cdot w}{\|w\|^2}.
\]
For the future-time problem \(t\ge0\), the minimizing time is
\(\max\{0,t_*\}\). The corresponding distance is obtained by
substituting that time in \(\|d+tw\|\).

Subtracting \(b+tv\) from both positions makes the second ship
stationary and preserves their relative distance at every time.
This invariance, rather than the mere existence of an easier
stationary-ship example, justifies the transfer from that special case.
\end{hbtheorem}

\section{Check the result against the original request}

P\'olya also uses concrete realizations of abstractions: a classroom
for the box, meeting room edges for coordinate axes, and a curve on
a rotating door for a swept surface of revolution. These examples
help specify or visualize a model; they do not add unmeasured
geometric facts to it. His mention of unusual members, such as a
whale within the class of mammals, illustrates the same logical
point as the specialization theorem: a universal condition has to
hold at every member of its declared domain.

\begin{hbmethod}{def:p12-heuristic-return}{A completed reasoning record}
Return the proposed answer or theorem together with its original
domain and conditions, the chain of justified transformations, the
resolution of introduced side conditions, and a forward verification.
For a request for all solutions, also justify that no cases were lost.
For a construction, verify that every operation was allowed.

Useful checks include units, symmetry, signs, limiting cases and a
second derivation. A failed check can reject an answer. Passing a
collection of checks establishes only their stated conditions unless
one has also proved that those conditions characterize the answer.
For example, dividing \(x(x-1)=0\) by \(x\) without separating
\(x=0\) loses a solution, and squaring the square-root equation
without its sign guard admits a false candidate.

Heuristics propose where to look. The formal implications and
reconstruction arguments determine what a successful step establishes.
No particular suggestion here has a proved bound on search effort
or a guarantee of finding a solution to every problem. An unresolved
search is not an impossibility proof, and an appealing analogy is not
an empirical guarantee of transfer.
\index{working backwards}
\index{heuristic reasoning!verification of transformations}
\end{hbmethod}
